\section*{Висновки}

\subsection*{Узагальнення результатів дослідження}

Проведене дослідження підтвердило гіпотезу про можливість радикального підвищення ефективності адміністративних процесів у закладах вищої освіти шляхом інтеграції генеративного штучного інтелекту. Реалізація переходу від реактивної звітності до проактивного управління науковою діяльністю є не лише технічно можливою, а й економічно доцільною навіть для ЗВО з обмеженими ресурсами.

Основні результати дослідження:
\begin{enumerate}[leftmargin=*,itemsep=0.25em]
\item \textbf{Розроблено концептуальну модель інтеграції GenAI} в адміністративні процеси ЗВО, що базується на принципі автономного агента (Autonomous Agent) з використанням механізму MCP Function Calling.
\item \textbf{Розроблено алгоритм мультиджерельного злиття даних «Golden Record»}, який забезпечує формування єдиної авторитетної версії запису публікації.
\item \textbf{Практична апробація} на прикладі системи ResearchIQ (15 дослідників, 3 кафедри) продемонструвала суттєве скорочення часу підготовки звітності та зниження частки дублікатів.
\item \textbf{Систематизовано світовий досвід} впровадження ШІ в управлінні дослідженнями та виділено критичні фактори успіху.
\item \textbf{Ідентифіковано ризики та обмеження} (галюцинації, bias, регуляторні вимоги) і запропоновано стратегії мітигації.
\end{enumerate}

\subsection*{Теоретичне значення}

Дослідження робить внесок у розвиток теорії цифрової трансформації публічного адміністрування, конкретизуючи роль GenAI як технологічного enabler проактивності та пропонуючи патерн MCP Function Calling Agent.

\subsection*{Практичне значення}

Результати дослідження мають безпосереднє практичне застосування для українських ЗВО, наукових підрозділів та регуляторів (policy-makers).

\subsection*{Дорожня карта впровадження}

Запропоновано трьохетапну дорожню карту: \textbf{Фундамент} (6--12 міс.), \textbf{Масштабування} (12--24 міс.), \textbf{Оптимізація} (24--36 міс.).

\subsection*{Обмеження дослідження}

Обмеження включають масштаб вибірки (15 дослідників), фокус на бібліометричній звітності та динаміку розвитку GenAI й регуляторного середовища.

\subsection*{Перспективи подальших досліджень}

Перспективи: рекомендаційні системи для вибору журналів, інтеграція з документообігом, етичні фреймворки для автономних агентів, порівняльний аналіз впроваджень, довгострокові ефекти автоматизації, мультимодальні можливості.

\bigskip

Підсумовуючи, інтеграція GenAI у систему адміністрування науки ЗВО є реалістичним напрямком цифрової трансформації, що може бути реалізований уже сьогодні за умови поєднання технологічних інновацій з організаційними змінами та дотримання принципу Human-in-the-Loop для критичних рішень.